\section{Background and Motivation}
\label{sec:background}

\subsection{Timed Automata and Industrial Networks}

Timed automata provide a standard formal model for systems whose
correctness depends on both discrete control decisions and quantitative
time constraints~\cite{alur1994theory,bengtsson2004timed}. Let $C$ be
a finite set of clocks. A clock constraint over $C$ is a conjunction of
atomic constraints of the form $x \sim b$, where $x \in C$,
$\sim \in \{<,\leq,=,\geq,>\}$, and $b$ is a non-negative numerical
bound. We write $\mathcal{B}(C)$ for the set of such constraints. A
timed automaton is a tuple
\[
  T=(L,\ell_0,C,\Sigma,\Theta,I),
\]
where $L$ is a finite set of locations, $\ell_0 \in L$ is the initial
location, $\Sigma$ is a set of action labels, $\Theta$ is a finite set
of transitions, and $I:L\rightarrow \mathcal{B}(C)$ assigns an invariant
to each location. A transition $\theta \in \Theta$ is written as
\[
  \theta=(\ell,g,a,R,\ell'),
\]
where $\ell$ and $\ell'$ are the source and target locations,
$g\in\mathcal{B}(C)$ is the guard, $a\in\Sigma$ is the action label,
and $R\subseteq C$ is the set of clocks reset by the transition.

The semantics is defined over states $(\ell,u)$, where $\ell$ is a
location and $u:C\rightarrow \mathbb{R}_{\geq 0}$ is a clock valuation.
A delay step lets time elapse in the current location as long as the
location invariant remains satisfied. A discrete step can be taken when
the current clock valuation satisfies the transition guard; the step
moves the automaton to the target location and resets the clocks in
$R$. Thus, guards constrain when discrete actions are enabled, whereas
invariants constrain how long the automaton may remain in a location.

Industrial timed-automata models are usually networks rather than single
automata. A network
\[
  \mathcal{N}=T_1 \parallel \cdots \parallel T_k
\]
composes templates for controllers, plants, environments, communication
protocols, schedulers, and monitors~\cite{larsen1997uppaal}. The exact
product semantics depends on the modeling language and tool, for
example on binary synchronization, broadcast channels, urgent or
committed locations, and data variables. In this work these structural
and data-level features are treated as fixed modeling context. The
repair space is deliberately restricted to numerical clock-bound
occurrences in guards and invariants; synchronization labels, reset
sets, clock references, locations, transitions, data updates, and
property formulas are not repair variables.

This restriction matches a common class of industrial modeling faults.
Timing constants often represent design thresholds, sampling periods,
timeout values, alarm durations, refractory intervals, or scheduling
deadlines. For example, vehicle controllers may need to complete a
gear-shifting sequence within a bounded response time
~\cite{lindahl1998gear}; railway models may encode gate, alarm, or
communication timeout requirements~\cite{basile2021unisig}; medical
devices may encode physiological timing windows~\cite{jiang2012pacemaker};
and avionics or aerospace models may encode execution-time bounds and
deadlines~\cite{han2018avionics}. In such models, an incorrect numerical
bound can make an otherwise intended behavior occur too early, too late,
or for an invalid duration.

We do not assume that all modeling errors are clock-bound errors.
Industrial models may also contain wrong resets, missing transitions,
incorrect synchronizations, erroneous data updates, or incomplete
properties. These faults are outside the current repair space. The
assumption in this paper is narrower: the discrete model structure is
treated as the intended design, while one or more numerical timing
bounds may be inaccurate due to calibration, transcription,
dimensioning, or design-space uncertainty. This is the fault model
targeted by clock-bound repair~\cite{koelbl2019clock} and is consistent
with mutation-based repair and testing, where faults are modeled as
small deviations from an otherwise near-correct artifact~\cite{demillo1978hints}.

\subsection{Properties and Timed Diagnostic Traces}

Timed-automata models are checked against timing properties using
real-time model checkers such as UPPAAL~\cite{larsen1997uppaal}. This
paper focuses on safety-style verification obligations. Direct examples
include queries of the form $A[]\,\psi$, where every reachable state must
satisfy the state predicate $\psi$, and deadlock-freedom checks such as
$A[]\,\mathit{not}\ \mathit{deadlock}$. Bounded-response and deadline
requirements are handled when they are encoded by observer or monitor
templates into safety-style violations, for example by reaching an
error or late location~\cite{vogel2022patterns}. Therefore, the repair
procedure reasons about property violations that can be represented as
reachable safety violations in the instrumented timed-automata network.

A discrete trace skeleton is a finite sequence of transitions
\[
  \tau=\theta_0,\theta_1,\ldots,\theta_{n-1}.
\]
A timed realization of $\tau$ is a sequence of non-negative delays and
clock valuations that makes the transition sequence executable from the
initial state while respecting guards, resets, invariants, and network
synchronization. The skeleton $\tau$ is feasible if it has at least one
timed realization. Given a property $\varphi$, a feasible skeleton is a
timed diagnostic trace (TDT) for $\varphi$ if at least one of its timed
realizations reaches a state that violates the safety condition
associated with $\varphi$~\cite{koelbl2019clock}. Depending on the
model checker, the reported counterexample may contain concrete delays,
symbolic zones, or only enough information to reconstruct the discrete
path and timing constraints.

A TDT is evidence of failure, not a root-cause explanation. It shows that
some timing realization of some discrete behavior violates a property,
but it does not identify which numerical bound is wrong, whether the
fault lies in a guard or invariant, or how the model should be changed.
Several bounds may jointly enable the violating realization, and one
incorrect bound may contribute to several diagnostic traces. Conversely,
blocking the discrete path of a TDT is usually not a valid repair: it
may remove intended functional behavior, such as a request,
acknowledgement, actuation, communication step, or scheduling decision.
A useful repair should change the timing envelope of the behavior while
preserving the relevant untimed functionality of the model.

\subsection{Clock-Bound Repair and Admissibility}

A repairable site is an occurrence of a numerical bound $b$ in an atomic
clock constraint $x\sim b$ that appears in a transition guard or a
location invariant. A clock-bound repair changes this occurrence to a
new value $b'$, equivalently by introducing a variation variable $v$ and
using the repaired bound $b+v$. The comparison operator, clock
reference, transition structure, location structure, synchronization
label, and reset set are kept fixed. Bounds are constrained to remain in
the numerical domain supported by the target model representation.

TDT-based clock-bound repair encodes the feasibility of a diagnostic
trace as constraints over delays and clock values, then introduces
variation variables for the bounds that occur in the encoded guards and
invariants~\cite{koelbl2019clock}. MaxSMT solving can then be used to
prefer small repairs, for example by maximizing soft constraints of the
form $v=0$ and, when needed, minimizing the magnitude of nonzero
variations~\cite{sebastiani2015optimization}. In the single-trace
setting, the resulting candidate is local: the same discrete trace
skeleton should remain feasible, but no feasible realization of that
skeleton should still witness the considered property violation.

Local trace repair is not sufficient for model repair. A bound change
that eliminates one violating realization may create a new violation on
another execution, fail another property, or make intended functionality
unavailable. MPTA-Repair therefore treats MaxSMT-based trace repair as
candidate generation only. A candidate is accepted only after two global
checks. First, the repaired model must satisfy the complete property set
under regression verification. Second, it must pass an admissibility
check based on TARTAR-style functional equivalence~\cite{koelbl2019clock}.

The admissibility check compares untimed functional behavior rather than
timed language equality. This distinction is necessary because a
successful timing repair is expected to remove some violating timed
realizations. Let $\Sigma_f$ denote the functional alphabet used for
admissibility, after excluding monitor-only, observer-only, and internal
labels when appropriate. The intended criterion is
\[
  \mathrm{Untimed}_{\Sigma_f}(M_{\mathit{bad}})
  =
  \mathrm{Untimed}_{\Sigma_f}(M_{\mathit{rep}}),
\]
meaning that the repaired model should not add or remove functional
action traces, although the timing of those traces may change. This
criterion prevents repairs that satisfy safety properties merely by
disabling intended behavior.

\subsection{Why Multi-Property Repair Is Needed}

Industrial timed-automata models are normally validated by a set of
properties rather than by a single assertion~\cite{vogel2022patterns}.
The properties jointly describe the reliability envelope of the model,
including absence of unsafe states, bounded response, timeout handling,
deadline satisfaction, admissible residence times, and deadlock freedom.
Let $M_{\mathit{bad}}$ be a faulty timed-automata network and let
\[
  \Phi=\{\varphi_1,\ldots,\varphi_m\}
\]
be the property set used to validate it. The goal is to construct a
repaired model $M_{\mathit{rep}}$ by changing a small number of
clock-bound constants in guards and invariants such that:
\[
  M_{\mathit{rep}} \models \Phi,
\]
the repaired bounds are well formed, and the repaired model preserves
the projected untimed functional behavior of $M_{\mathit{bad}}$.

The need for joint reasoning can be seen in the gear-shift controller
used as a running example. The driver periodically issues a shift
request, and the controller performs torque cut, gear engagement,
acknowledgement, and recovery. Suppose the faulty controller uses
bounds such as $\mathit{cut}\leq 1$, $\mathit{cut}\geq 1$,
$\mathit{mesh}\leq 3$, and $\mathit{mesh}\geq 3$. One property may
require at least two time units of torque cut before gear engagement,
while another property may require acknowledgement within four time
units of the request. A single-property repair that changes only the
$\mathit{cut}$ bounds from $1$ to $2$ can block the early-engagement
violation. However, if the $\mathit{mesh}$ bounds remain at $3$, the
earliest acknowledgement may occur at time $5$, thereby violating the
response deadline. A multi-property repair must instead consider both
requirements and may need to change both the $\mathit{cut}$ and
$\mathit{mesh}$ timing bounds while preserving the same untimed
request--engage--acknowledge behavior.

This example illustrates why violated properties, diagnostic traces, and
repair variables should not be treated as unrelated subproblems.
Different properties may share clocks, locations, synchronizations, or
execution fragments. Multiple TDTs may expose different realizations of
the same underlying timing defect. Conversely, a repair computed from
one trace can be locally valid yet globally unacceptable after full
regression. MPTA-Repair addresses this setting by accumulating
multi-property and multi-counterexample repair evidence, generating
small clock-bound candidates with MaxSMT, and accepting a candidate only
after complete property regression and functional-admissibility
checking.